\originalchapter{课程模拟考试与中文题解}
本部分题面来自 MIT 18.703, Spring 2013 的3份公开模拟试卷, 原编号、分值及小问保持对应. 官网没有随卷提供答案, 以下解答均为本整理本独立撰写并交叉审读, 不是官方答案. 原卷中的空白答题页合并, 不影印英文整页. 正式两次期中及期末试卷未在此课程考试页面公开.
\originalsection{第一次模拟期中}
原卷限时80分钟, 闭卷、不可使用计算器, 共7题; 题目95分, 表达5分, 合计100分.
\begin{exercise}\label{ex:m1q1}原卷第1题 (15分). (i) 定义正规子群. (ii) 定义群同态. (iii) 定义交错群 $A_n$.\end{exercise}
\begin{solution}
(i) $N\le G$ 且任意 $g$ 有 $gNg^{-1}=N$. (ii) $\varphi:G\to H$ 满足 $\varphi(ab)=\varphi(a)\varphi(b)$. (iii) $A_n$ 是 $S_n$ 中偶置换组成的子群, 即符号同态的核; 偶性由换位分解的偶数长度定义, 其不依赖分解已在第3章证明.
\end{solution}
\begin{exercise}原卷第2题 (15分). (i) 给 $A_4$ 的一个非平凡真正规子群 $V$, 说明其同构类型. (ii) 写其全部左陪集. (iii) 求 $A_4/V$ 的同构类型.\end{exercise}
\begin{solution}
取 $V=\{e,(12)(34),(13)(24),(14)(23)\}\cong C_2\times C_2$, 共轭保留双换位, 故正规. 全部左陪集为
\begin{align*}
V,\quad &(123)V=\{(123),(134),(243),(142)\},\\
&(132)V=\{(132),(234),(124),(143)\}.
\end{align*}
各集合含4个元素且不交, 共覆盖12个元素. 商阶3, 所以为 $C_3$.
\end{solution}
\begin{exercise}原卷第3题 (20分). 对 $H\le G$, 证明下列条件等价: (1) $H$ 正规; (2) 每个 $g$ 有 $gHg^{-1}=H$; (3) 每个 $a$ 有 $aH=Ha$; (4) 全部左陪集组成的集合与全部右陪集组成的集合相同.\end{exercise}
\begin{solution}
(1)与(2)为定义, (2)右乘 $g$ 等价于(3), (3)推出(4). 若(4)成立, 对任一 $a$, 存在 $b$ 使 $aH=Hb$. 因 $a\in Hb$, 有 $a=hb$, $h\in H$, 故 $Ha=Hhb=Hb=aH$, 得(3). 最后一步不能仅从两边“陪集数相同”推出, 必须用到同一陪集中包含 $a$.
\end{solution}
\begin{exercise}原卷第4题 (10分). 设 $N\trianglelefteq G$. 证明 $G/N$ 交换当且仅当 $N$ 包含所有交换子 $aba^{-1}b^{-1}$.\end{exercise}
\begin{solution}
商中 $aN,bN$ 交换当且仅当其交换子陪集为 $N$, 即 $aba^{-1}b^{-1}\in N$. 对所有 $a,b$ 量化, 即得两方向.
\end{solution}
\begin{exercise}\label{ex:commute}原卷第5题 (10分), 与作业5第12题同题. 若 $H,K\trianglelefteq G$ 且交为单位子群, 证明其元素两两交换.\end{exercise}
\begin{solution}
对 $h\in H,k\in K$, $hkh^{-1}k^{-1}$ 由 $K$ 正规属于 $K$, 由 $H$ 正规又属于 $H$. 交平凡使该交换子为 $e$, 即 $hk=kh$.
\end{solution}
\begin{exercise}\label{ex:sylowstate}原卷第6题 (15分). (i) 陈述 Sylow 定理. (ii) 证明不同素数 $p,q$ 的乘积阶群不为单群.\end{exercise}
\begin{solution}
(i) 若 $|G|=p^am$, $p\nmid m$, 则有阶 $p^a$ 子群; 所有这样的子群共轭, 任意 $p$ 子群包含在其中一个内; 个数整除 $m$ 且模 $p$ 余1. 完整证明见定理\ref{thm:sylow}. (ii) 不妨 $p<q$, $n_q\mid p$, $n_q\equiv1\pmod q$, 所以 $n_q=1$. 唯一 Sylow $q$ 子群非平凡、真且正规, 排除单性.
\end{solution}
\begin{exercise}原卷第7题 (10分). 建立包含 $N\trianglelefteq G$ 的子群与 $G/N$ 子群的自然双射, 并证明正规性保持.\end{exercise}
\begin{solution}
对投影 $\pi$, 用 $H\mapsto\pi(H)$ 与 $\bar H\mapsto\pi^{-1}(\bar H)$. 若 $N\subseteq H$, 则 $\pi^{-1}(\pi(H))=H$, 因像相同的两元素相差核元素; 满射性给反向复合为恒等. 子群封闭性由像、原像保持乘法和逆元得出. $\pi(gHg^{-1})=\pi(g)\pi(H)\pi(g)^{-1}$ 且原像唯一, 故正规性双向对应.
\end{solution}
\originalsection{第二次模拟期中}
原卷限时80分钟, 闭卷、不可使用计算器, 共6题; 题目95分, 表达5分, 合计100分.
\begin{exercise}原卷第1题 (15分). (i) 定义整环的不可约元. (ii) 定义素元. (iii) 定义主理想整环.\end{exercise}
\begin{solution}
不可约元是非零非单位, 任一因子分解都含单位因子. 素元是非零非单位, 若整除乘积则整除至少一个因子. PID 是每个理想均由一个元素生成的整环. “整环”及“非单位”条件均不能省略.
\end{solution}
\begin{exercise}\label{ex:120}原卷第2题 (15分). (i) 陈述 Sylow 定理. (ii) 证明不存在阶120的单群.\end{exercise}
\begin{solution}
(i) 见习题\ref{ex:sylowstate}的完整陈述. (ii) 设 $G$ 单, 则 $n_5\mid24$, $n_5\equiv1\pmod5$, 只能为1或6, 单性排除1. 共轭作用在六个 Sylow 5子群上给非平凡同态 $G\to S_6$, 核正规故单射. 符号同态在 $G$ 上也必平凡 (否则单射到阶2群), 所以像在 $A_6$. 于是 $G$ 在 $A_6$ 中指数3. $A_6$ 在其三个左陪集上作用, 得非平凡 $A_6\to S_3$. 作业6第12题的证明给 $A_6$ 单, 迫使此映射单射, 但 $360>6$, 矛盾.
\end{solution}
\begin{exercise}原卷第3题 (15分). 对整环 $R$ 的真理想 $I$, 证明 $R/I$ 是域当且仅当 $I$ 极大.\end{exercise}
\begin{solution}
事实上只需 $R$ 交换有1. 若 $I$ 极大, 对 $a\notin I$, $I+(a)=R$, 所以 $i+ra=1$ 给商中逆元. 若商为域, 任意 $J\supsetneq I$ 中的非零像可逆, 故 $J/I$ 含1, 于是 $J=R$. 真理想条件排除零商环.
\end{solution}
\begin{exercise}原卷第4题 (15分). (i) 证明理想升链的并是理想. (ii) 证明 PID 的理想升链最终稳定.\end{exercise}
\begin{solution}
若 $x,y$ 属并, 它们同属较后一个理想, 故差也在并; 任意环元素乘 $x$ 仍在含 $x$ 的理想. 若并 $I=(d)$, 元素 $d$ 已在某个 $I_N$, 故 $I\subseteq I_N$, 而反向包含本来成立. 所有后续理想都等于 $I_N$.
\end{solution}
\begin{exercise}原卷第5题 (20分). (i) 证明两理想的交是理想. (ii) 并集必为理想吗?\end{exercise}
\begin{solution}
交中的差以及与任意环元素的左右乘积同时在两个理想中, 故在交中. 并不必是理想, 例如 $2\Z\cup3\Z$ 含2和3却不含5. 一般并是加法子群当且仅当一个理想包含另一个: 若各取一个不属于对方的元素, 它们的和不属于任何一方.
\end{solution}
\begin{exercise}原卷第6题 (15分). 在 PID 中证明非零 $a,b$ 有 gcd $d$ 且 $d=ra+sb$.\end{exercise}
\begin{solution}
写理想 $(a,b)=(d)$. 因 $a,b\in(d)$, $d$ 整除二者. 生成元 $d$ 属 $(a,b)$ 故具有所述表达, 任意公因子由表达式整除 $d$. 这就是 gcd 定义, 唯一性仅差单位.
\end{solution}
\originalsection{模拟期末}
原卷限时3小时, 闭卷、不可使用计算器, 共11题; 题目175分, 表达5分, 合计180分.
\begin{exercise}原卷第1题 (30分). 定义 (i) 群; (ii) 群自同构; (iii) 二面体群; (iv) 理想; (v) PID; (vi) UFD.\end{exercise}
\begin{solution}
(i) 群的运算结合, 有双侧单位元及每个元素的双侧逆元. (ii) 自同构是群到自身的双射同态. (iii) $D_n$ 为正 $n$ 边形的全对称群, 本书阶约定为 $2n$; 可用 $r^n=s^2=e,srs=r^{-1}$ 描述. (iv) 理想是环的加法子群, 对任意环元素左右乘法吸收; 交换环只需一侧. (v) PID 为每个理想主的整环. (vi) UFD 为每个非零非单位都可分解成不可约元且仅差次序与相伴的整环.
\end{solution}
\begin{exercise}原卷第2题 (15分). (i) 证明共轭关系是等价关系. (ii) 列出 $S_5$ 的等价类.\end{exercise}
\begin{solution}
用单位元给自反, 共轭元的逆给对称, 两次共轭用乘积给传递. $S_5$ 的共轭类按5的分拆列出, 类型为 $1^5,2\,1^3,2^2\,1,3\,1^2,3\,2,4\,1,5$. 其大小分别为 $1,10,15,20,20,30,24$, 和为120. 大小公式 $5!/\prod_r(r^{m_r}m_r!)$ 来自循环内起点 $r$ 种重复与相同长度循环的 $m_r!$ 种重排.
\end{solution}
\begin{exercise}\label{ex:smallgroups}原卷第3题 (15分). 分类所有阶至多10的群.\end{exercise}
\begin{solution}
各阶的全部同构类型如下: 1阶平凡群; 2,3,5,7阶分别唯一 $C_2,C_3,C_5,C_7$; 4阶 $C_4,C_2^2$; 6阶 $C_6,S_3$; 8阶 $C_8,C_4\times C_2,C_2^3,D_4,Q_8$; 9阶 $C_9,C_3^2$; 10阶 $C_{10},D_5$.

素数阶与平方阶已由拉格朗日及第4章给出. 对 $2q$ ($q=3,5$), 唯一 Sylow $q$ 子群 $\langle a\rangle$ 正规, 柯西定理给阶2元素 $b$ 在其外. 群中每个元素为 $a^i$ 或 $a^ib$, 且 $bab^{-1}=a^u$, $u^2=1$ 模 $q$, 故 $u=1$ 或 $-1$. 前者循环直积, 后者二面体; $q=3$ 的二面体即 $S_3$.

8阶交换群由有限生成交换群结构定理给三类. 若非交换, 不能有阶8元素. 若所有非单位元阶2, 由 $(ab)^2=e$ 推得 $ab=ba$, 矛盾. 故取阶4元素 $a$, 子群指数2而正规, 选 $b$ 在其外. 共轭作用非平凡, 必 $bab^{-1}=a^{-1}$; 否则 $a,b$ 交换使群交换. $b^2\in\langle a\rangle$, 又与 $b$ 交换, 因此 $b^2=a^k$ 的 $k$ 必为0或2. $b^2=e$ 得 $D_4$, $b^2=a^2$ 得四元数群 $Q_8=\{\pm1,\pm i,\pm j,\pm k\}$, 其中 $i^2=j^2=k^2=ijk=-1$. 两个表示各最多8个规范词, 具体群已有8个元素, 因而无遗漏. $D_4$ 有5个阶2元素, $Q_8$ 只有1个, 故不同构.
\end{solution}
\begin{exercise}原卷第4题 (15分). (i) 陈述第二同构定理. (ii) 证明它.\end{exercise}
\begin{solution}
若 $H\le G,N\trianglelefteq G$, 则 $H/(H\cap N)\cong HN/N$. $HN$ 由正规性封闭; $H\to HN/N$, $h\mapsto hN$, 满射且核为 $H\cap N$, 第一同构定理给结论. 原课程对“第二”编号采用此版本, 不与第三同构定理混淆.
\end{solution}
\begin{exercise}原卷第5题 (15分). (i) 陈述 Sylow 定理. (ii) 若 $|G|=pqr$, 三素数互异, 证明 $G$ 不单.\end{exercise}
\begin{solution}
(i) 见习题\ref{ex:sylowstate}. (ii) 设 $p<q<r$ 且群单. $n_r\mid pq$, 模 $r$ 余1且大于1, 不可能为 $p$ 或 $q$, 只能 $pq$. 各阶 $r$ 子群除单位外不交, 共 $pq(r-1)$ 个元素. 又 $n_q$ 是 $pr$ 的因子, 大于1且模 $q$ 余1, 不可能为 $p$, 所以 $n_q\ge r$. 它们至少给 $r(q-1)$ 个不同阶 $q$ 元素, 与前者不交. 但
\[
pq(r-1)+r(q-1)=pqr+\bigl(r(q-1)-pq\bigr)>pqr,
\]
因 $r>q$ 且 $p\le q-1$ 得 $r(q-1)>q(q-1)\ge pq$. 矛盾.
\end{solution}
\begin{exercise}原卷第6题 (15分). 设 $R$ 交换, $I,J$ 为理想. (i) 素理想 $P\supseteq IJ$ 时证明 $P\supseteq I$ 或 $P\supseteq J$. (ii) 建立 $R/IJ$ 与 $R/(I\cap J)$ 的素理想自然双射. (iii) 给 $IJ\ne I\cap J$ 的例子.\end{exercise}
\begin{solution}
(i) 若 $I\not\subseteq P$, 取 $i\in I\setminus P$, 则对所有 $j\in J$, $ij\in P$, 素性迫使 $j\in P$. (ii) $IJ\subseteq I\cap J$. 包含 $IJ$ 的素理想由(i)包含一方, 因而包含交; 包含交当然包含积. 两商环的素理想因此对应同一批 $R$ 的素理想, 用投影原像与像得到双射. (iii) $R=\Z,I=J=2\Z$, 则积 $4\Z$, 交 $2\Z$.
\end{solution}
\begin{exercise}原卷第7题 (10分). 非域 UFD 是否必有无限多个两两不相伴的不可约元?\end{exercise}
\begin{solution}
否. 取 $R=\{a/b\in\Q:b\text{ 为奇数}\}=\Z_{(2)}$. 非零元素唯一写为 $2^ku$, $k\ge0$, $u$ 为单位 (分子分母皆奇). 非单位正是 $k>0$, 不可约元正是 $k=1$, 全与2相伴. 分解存在且唯一, 所以为非域 UFD, 却只有一个不可约元相伴类.
\end{solution}
\begin{exercise}原卷第8题 (10分). 举一个每个非零非单位可分解为不可约元却非 UFD 的整环.\end{exercise}
\begin{solution}
取 $\Z[\sqrt{-5}]$. 范数为正整数且每个非单位范数大于1, 严格下降归纳保证有限分解. $6=2\cdot3=(1+\sqrt{-5})(1-\sqrt{-5})$ 给不相伴的两种不可约分解; 具体不可约性见第6章的完整范数检验.
\end{solution}
\begin{exercise}原卷第9题 (15分). (i) 证明 $\Z[i]$ 欧几里得. (ii) $6-i$ 是否为素元?\end{exercise}
\begin{solution}
(i) 用范数 $a^2+b^2$; 对 $z/w$ 实虚部各选最近整数 $q$, 余项 $r=z-qw$ 满足 $N(r)\le N(w)/2<N(w)$. (ii) 范数 $37$ 为有理素数. 若 $6-i=ab$, 范数乘法性迫使一个因子范数1, 即单位, 故不可约. 欧几里得整环是 PID, 不可约即素元, 所以是.
\end{solution}
\begin{exercise}原卷第10题 (10分). 写出 $\F_5[x]$ 的全部不可约二次多项式.\end{exercise}
\begin{solution}
先列首一式. 对 $x^2+bx+c$, 完全平方表明无根当且仅当判别式 $b^2-4c$ 为非平方, 即2或3. 全部十个首一式为
\begin{align*}
&x^2+2,\ x^2+3,\ x^2+x+1,\ x^2+x+2,\\
&x^2+2x+3,\ x^2+2x+4,\ x^2+3x+3,\ x^2+3x+4,\\
&x^2+4x+1,\ x^2+4x+2.
\end{align*}
原题未限定首一, 所以还须取上述每式的任一非零常数倍, 共40个二次不可约多项式.
\end{solution}
\begin{exercise}原卷第11题 (25分). (i) 陈述高斯引理与艾森斯坦判别. (ii) 证明 $1+x^3+x^6$ 在 $\Q[x]$ 不可约. (iii) 证明 $1-t^2+t^5$ 在 $\Q[t]$ 不可约.\end{exercise}
\begin{solution}
(i) 本原多项式之积本原, 本原整数多项式在 $\Z[x]$ 与 $\Q[x]$ 的不可约性等价. 艾森斯坦条件为: 正次数整数多项式存在素数 $p$ 整除全部非最高次系数, $p$ 不整除最高系数, 且 $p^2$ 不整除常数项; 则该多项式在 $\Q[x]$ 上不可约. 完整证明见引理\ref{lem:gauss}及定理\ref{thm:eisenstein}. (ii) 令 $x=y+1$, 得
\[
y^6+6y^5+15y^4+21y^3+18y^2+9y+3.
\]
素数3整除所有非最高系数, 不整除最高系数且9不整除常数3, 判别成立; 平移可逆, 故原式不可约.

(iii) 模2为 $h=t^5+t^2+1$, 在0,1无根. 若五次式可约且无一次因子, 必有不可约二次因子, 因可能次数分拆至少一项不大于2. $\F_2$ 唯一不可约二次为 $t^2+t+1$. 模此式 $t^2=t+1$, $t^3=1$, $t^5=t^2$, 故 $h\equiv1$, 不整除. 因而 $h$ 不可约, 约化判别给有理不可约性.
\end{solution}
